\chapter{\code{StringInternerBackend} Trait}
\label{ch:interner-trait}

\begin{note}
Stub. See Chapter~\ref{ch:interner} for the underlying struct and
metrics. The trait exists as a seam for a future
\code{GenerationalInterner} or \code{PartitionedInterner}; no
second impl is shipped today, the trait is pre-positioned for when
production evidence justifies one.
\end{note}

\section{Contract}

\begin{lstlisting}[style=rust]
pub trait StringInternerBackend {
    fn intern(&mut self, s: &str) -> StrId;
    fn resolve(&self, id: StrId) -> &str;
    fn try_resolve(&self, id: StrId) -> Option<&str>;
    fn get(&self, s: &str) -> Option<StrId>;
    fn len(&self) -> usize;
    fn is_empty(&self) -> bool;
    fn reserve(&mut self, additional: usize);
    fn metrics(&self) -> InternerMetrics;
}

pub struct InternerMetrics {
    pub unique_strings: usize,
    pub approx_bytes:   usize,
}
\end{lstlisting}

\section{Trigger conditions}

A second backend is worth implementing when any of:
\begin{itemize}
    \item \code{metrics().unique\_strings} exceeds $\sim 10^7$ in
        a live session.
    \item A streaming session runs 24$\times$7 for more than 72h
        without compaction.
    \item A single process serves multiple tenants
        (MSSP deployment).
\end{itemize}

The engine exposes the metrics via
\code{api.get\_interner\_metrics()} so an operator can observe the
approach to a trigger rather than hit it blind.

\begin{inthecode}
\filepath{graph\_hunter\_core/src/interner.rs} --- trait + metrics.\\
\filepath{TODO.md} --- trigger criteria.
\end{inthecode}
